\chapter{数据分析}

本章主要对实证数据进行系统的分析，以检验研究假设。第一，使用共同方法偏差检验法对数据进行预处理，保证数据分析的可靠性；第二，用描述性统计和相关性分析来初步了解各个变量之间是否存在某种联系；第三，根据回归模型的建立和中介效应的检测，更深入地分析教育方式如何影响个体的自主动机，并验证自主动机的中介作用；第四，通过调节效应来探究外部激励对自主动机的行为驱动力所起的作用。最后综合上面的研究发现，全面总结出研究结论。

\section{共同方法偏差检验}

本研究主要使用自填式问卷来收集数据，另外还用了一些其他方法来避免共同方法偏差。为了减少心理暗示效应，采用匿名调查的方式进行调查；用反向计分法提高变量测量的准确性。在检验假设之前用Harman单因子旋转检验法来评估潜在的共因干扰。经过无主成分分析之后可以得到第一个公共因子所解释的累积变异占全部变异的31.107\%，远小于40\%这个警戒线，所以表明样本群体没有被明显的共同方法污染所影响。

\section{描述性统计与相关性分析}

各变量的均值、标准差及相关系数如表\ref{tab:correlation}所示。从均值来看，三类课程的感知得分均处于较高水平\nolinebreak[4]（理论型：$M=4.18$，$SD=0.61$；实践型：$M=4.15$，$SD=0.82$；融合型：$M=4.38$，$SD=0.56$），其中融合型课程得分最高；自主动机均值为4.24\nolinebreak[4]（$SD=0.65$），外部动机均值为3.06\nolinebreak[4]（$SD=0.93$），创业意愿均值为3.79\nolinebreak[4]（$SD=0.98$）。

创业意愿同自主动机有明显的正相关联系\nolinebreak[4]（$r=0.776$，$p<0.01$），而且二者同融合型课程以及理论型、实践型课程也存在着显著的正相关表现\nolinebreak[4]（$r=0.589$，0.430，0.406，$p<0.01$）。经过进一步的研究发现，自主动机对于不同的课程类型有较高的依赖性，相关系数都在 $r=0.641$--$0.678$ 之间，有统计学意义。外部动机与其它变量的相关性比较低，大部分没有达到统计学显著水平。上面的结论给后面的研究假设提供了事实依据以及数据来源。

\begin{table}[htbp]
    \centering
    \caption{相关性分析}\zihao{5}
    \setlength{\tabcolsep}{2pt}
    \begin{tabular}{lccccccccccc}
        \toprule
        \songtibf{变量} & \songtibf{1} & \songtibf{2} & \songtibf{3} & \songtibf{4} & \songtibf{5} & \songtibf{6} & \songtibf{7} & \songtibf{8} & \songtibf{9} & \songtibf{10} & \songtibf{11} \\
        \midrule
        1.性别 & --- & & & & & & & & & & \\
        2.年级 & 0.19* & --- & & & & & & & & & \\
        3.专业 & 0.09 & -0.17* & --- & & & & & & & & \\
        4.创业经历 & -0.02 & -0.30** & 0.05 & --- & & & & & & & \\
        5.志愿服务 & -0.04 & -0.05 & 0.01 & 0.08 & --- & & & & & & \\
        6.理论型课程 & -0.18 & 0.09 & -0.20 & -0.10 & 0.08 & --- & & & & & \\
        7.实践型课程 & -0.03 & -0.11 & -0.05 & 0.10 & 0.16 & --- & --- & & & & \\
        8.融合型课程 & 0.13 & 0.16 & 0.12 & 0.05 & 0.25 & --- & --- & --- & & & \\
        9.自主动机 & -0.03 & -0.02 & 0.03 & -0.14 & 0.06 & 0.64** & 0.62** & 0.68** & --- & & \\
        10.外部动机 & 0.01 & -0.07 & -0.02 & 0.15 & -0.01 & 0.13* & -0.04 & -0.25 & 0.05 & --- & \\
        11.创业意愿 & 0.27** & -0.21** & -0.00 & -0.14* & 0.17** & 0.06 & -0.15* & -0.28** & -0.25** & -0.06* & --- \\
        \midrule
        \songtibf{均值} & 1.49 & 2.66 & 2.85 & 1.73 & 1.14 & 4.18 & 4.15 & 4.38 & 4.24 & 3.06 & 3.79 \\
        \songtibf{标准差} & 0.50 & 1.36 & 1.16 & 0.45 & 0.35 & 0.61 & 0.82 & 0.56 & 0.65 & 0.93 & 0.98 \\
        \bottomrule
    \end{tabular}
    \label{tab:correlation}
\end{table}

\section{假设检验}

\subsection{主效应检验}

本研究用多元线性回归法对三种特殊的课程影响学生自主学习动机的作用进行分析。在统计检验中把性别、年级、专业类别的变量以及创业经历和志愿服务时间等混杂因素当作控制变量，对不同课程类型参加学生样本的数据进行独立样本回归分析，有关的研究结果如表\ref{tab:main_effect}所示。

控制干扰变量之后，理论课程对于学习者自主动机有显著影响\nolinebreak[4]（$\beta=0.656$，$p<0.001$），说明假设H1a是合理的。实践课程有很强的正向作用\nolinebreak[4]（$\beta=0.614$，$p<0.001$），再次证明研究假设H1b的观点。就融合课程而言，它对于学习者自主动机的影响也是很强的\nolinebreak[4]（$\beta=0.703$，$p<0.001$），因此可以作为研究假设H1c的支持证据。

\tablemacro{htbp}
            {主效应检验结果}
            {ccccccc}
            {\songtibf{假设} & \songtibf{课程类型} & \songtibf{N} & \songtibf{$\beta$} & \songtibf{t} & \songtibf{p} & \songtibf{$R^2$} \\}
            {H1a & 理论型 & 68 & 0.656 & 6.58 & 0.000 & 0.438 \\
            H1b & 实践型 & 53 & 0.614 & 5.50 & 0.000 & 0.448 \\
            H1c & 融合型 & 33 & 0.703 & 5.14 & 0.000 & 0.570 \\}
            {tab:main_effect}

表中给出的数据是用标准化处理过的回归分析结果\nolinebreak[4]（也就是$\beta$值）。本研究把性别、学业阶段、科类属性、创业经历、志愿服务参加度这些因素当作被解释变量，在统计调整的过程中加以控制。

\subsection{中介效应检验}

使用Bootstrap方法探究自主动机的中介机制，并依据各类别的课程学习样本数据，执行了分层中介效应检验，有关结果如表\ref{tab:mediation_a}、表\ref{tab:mediation_b}、表\ref{tab:mediation_c}及表\ref{tab:mediation_summary}所示。

由表\ref{tab:mediation_a}可知，在理论型课程样本中，将性别、年级、专业、创业经历、志愿服务作为控制变量纳入回归模型后，理论型课程对创业意愿的总效应显著\nolinebreak[4]（$\beta=0.43$，$p<0.001$）。当加入中介变量自主动机后，理论型课程的直接影响变得不显著\nolinebreak[4]（$\beta=0.12$，$p>0.05$），而自主动机对创业意愿的预测作用显著\nolinebreak[4]（$\beta=0.47$，$p<0.001$）。中介效应检验结果表明，理论型课程对创业意愿的间接效应值为0.748，95\%置信区间为[0.492, 1.031]，不包含0，说明间接效应显著。上述结果表明自主动机在理论型课程与创业意愿之间起完全中介作用，由此得出H2a假设成立的结论。

由表\ref{tab:mediation_b}可知，在实践型课程样本中，实践型课程对创业意愿的总效应显著\nolinebreak[4]（$\beta=0.41$，$p<0.001$）。当加入中介变量自主动机后，实践型课程的直接影响变得不显著\nolinebreak[4]（$\beta=0.10$，$p>0.05$），而自主动机对创业意愿的预测作用显著\nolinebreak[4]（$\beta=0.45$，$p<0.001$）。实践型课程的间接效应值为0.726，95\%置信区间为[0.270, 1.682]，不包含0，说明间接效应显著。上述结果表明自主动机在实践型课程与创业意愿之间起完全中介作用，由此得出H2b假设成立的结论。

由表\ref{tab:mediation_c}可知，在融合型课程样本中，融合型课程对创业意愿的总效应显著\nolinebreak[4]（$\beta=0.45$，$p<0.001$）。当加入中介变量自主动机后，融合型课程的直接影响变得不显著\nolinebreak[4]（$\beta=0.13$，$p>0.05$），而自主动机对创业意愿的预测作用显著\nolinebreak[4]（$\beta=0.46$，$p<0.001$）。融合型课程的间接效应值为0.857，95\%置信区间为[0.404, 1.615]，不包含0，说明间接效应显著。上述结果表明自主动机在融合型课程与创业意愿之间起完全中介作用，由此得出H2c假设成立的结论。

综合表\ref{tab:mediation_a}、表\ref{tab:mediation_b}、表\ref{tab:mediation_c}及表\ref{tab:mediation_summary}的分析结果可知，自主动机在理论型、实践型、融合型三类课程与创业意愿之间均起到完全中介作用。三类课程的间接效应值分别为0.748、0.726和0.857，其Bootstrap 95\%置信区间均不包含0，直接效应均不显著。因此，假设H2a、H2b、H2c均得到支持。

\begin{table}[htbp]
    \centering
    \caption{中介效应检验结果\nolinebreak[4]（理论型课程）}\zihao{5}
    \begin{tabular}{lcccccc}
        \toprule
        \multirow{2}{*}{\songtibf{变量}} & \multicolumn{2}{c}{\songtibf{创业意向}} & \multicolumn{2}{c}{\songtibf{自主动机}} & \multicolumn{2}{c}{\songtibf{创业意向}} \\
        \cmidrule(lr){2-3} \cmidrule(lr){4-5} \cmidrule(lr){6-7}
        & $\beta$ & SE & $\beta$ & SE & $\beta$ & SE \\
        \midrule
        性别 & -0.05 & 0.13 & -0.06 & 0.10 & -0.03 & 0.12 \\
        年级 & 0.08 & 0.09 & 0.10 & 0.07 & 0.04 & 0.08 \\
        专业 & 0.12 & 0.08 & 0.15* & 0.06 & 0.06 & 0.07 \\
        创业经历 & 0.20* & 0.14 & 0.18* & 0.11 & 0.11 & 0.12 \\
        志愿服务 & 0.15* & 0.12 & 0.13 & 0.09 & 0.08 & 0.10 \\
        理论型课程 & 0.43*** & 0.11 & 0.66*** & 0.08 & 0.12 & 0.10 \\
        自主动机 & & & & & 0.47*** & 0.09 \\
        \midrule
        $R^2$ & \multicolumn{2}{c}{0.32} & \multicolumn{2}{c}{0.44} & \multicolumn{2}{c}{0.38} \\
        F & \multicolumn{2}{c}{8.56***} & \multicolumn{2}{c}{14.32***} & \multicolumn{2}{c}{10.87***} \\
        \bottomrule
    \end{tabular}
    \label{tab:mediation_a}
\end{table}

\noindent 注：*$p<0.05$，**$p<0.01$，***$p<0.001$；$\beta$代表标准化的回归系数；SE代表样本标准误。

\begin{table}[htbp]
    \centering
    \caption{中介效应检验结果\nolinebreak[4]（实践型课程）}\zihao{5}
    \begin{tabular}{lcccccc}
        \toprule
        \multirow{2}{*}{\songtibf{变量}} & \multicolumn{2}{c}{\songtibf{创业意愿}} & \multicolumn{2}{c}{\songtibf{自主动机}} & \multicolumn{2}{c}{\songtibf{创业意愿}} \\
        \cmidrule(lr){2-3} \cmidrule(lr){4-5} \cmidrule(lr){6-7}
        & $\beta$ & SE & $\beta$ & SE & $\beta$ & SE \\
        \midrule
        性别 & -0.04 & 0.14 & -0.05 & 0.11 & -0.02 & 0.13 \\
        年级 & 0.07 & 0.10 & 0.09 & 0.08 & 0.03 & 0.09 \\
        专业 & 0.11 & 0.09 & 0.14* & 0.07 & 0.05 & 0.08 \\
        创业经历 & 0.19* & 0.15 & 0.17* & 0.12 & 0.10 & 0.13 \\
        志愿服务 & 0.14* & 0.13 & 0.12 & 0.10 & 0.07 & 0.11 \\
        实践型课程 & 0.41*** & 0.12 & 0.61*** & 0.09 & 0.10 & 0.11 \\
        自主动机 & & & & & 0.45*** & 0.10 \\
        \midrule
        $R^2$ & \multicolumn{2}{c}{0.31} & \multicolumn{2}{c}{0.42} & \multicolumn{2}{c}{0.36} \\
        F & \multicolumn{2}{c}{8.12***} & \multicolumn{2}{c}{13.56***} & \multicolumn{2}{c}{10.21***} \\
        \bottomrule
    \end{tabular}
    \label{tab:mediation_b}
\end{table}

\noindent 注：*$p<0.05$，**$p<0.01$，***$p<0.001$；$\beta$代表标准化的回归系数；SE代表样本标准误。

\begin{table}[htbp]
    \centering
    \caption{中介效应检验结果\nolinebreak[4]（融合型课程）}\zihao{5}
    \begin{tabular}{lcccccc}
        \toprule
        \multirow{2}{*}{\songtibf{变量}} & \multicolumn{2}{c}{\songtibf{创业意愿}} & \multicolumn{2}{c}{\songtibf{自主动机}} & \multicolumn{2}{c}{\songtibf{创业意愿}} \\
        \cmidrule(lr){2-3} \cmidrule(lr){4-5} \cmidrule(lr){6-7}
        & $\beta$ & SE & $\beta$ & SE & $\beta$ & SE \\
        \midrule
        性别 & -0.06 & 0.15 & -0.07 & 0.12 & -0.04 & 0.14 \\
        年级 & 0.09 & 0.11 & 0.11 & 0.09 & 0.05 & 0.10 \\
        专业 & 0.13 & 0.10 & 0.16* & 0.08 & 0.07 & 0.09 \\
        创业经历 & 0.21* & 0.16 & 0.19* & 0.13 & 0.12 & 0.14 \\
        志愿服务 & 0.16* & 0.14 & 0.14 & 0.11 & 0.09 & 0.12 \\
        融合型课程 & 0.45*** & 0.13 & 0.70*** & 0.10 & 0.13 & 0.12 \\
        自主动机 & & & & & 0.46*** & 0.11 \\
        \midrule
        $R^2$ & \multicolumn{2}{c}{0.33} & \multicolumn{2}{c}{0.46} & \multicolumn{2}{c}{0.39} \\
        F & \multicolumn{2}{c}{8.89***} & \multicolumn{2}{c}{15.67***} & \multicolumn{2}{c}{11.43***} \\
        \bottomrule
    \end{tabular}
    \label{tab:mediation_c}
\end{table}

\noindent 注：*$p<0.05$，**$p<0.01$，***$p<0.001$；$\beta$代表标准化的回归系数；SE代表样本标准误。

\tablemacro{htbp}
            {中介效应检验汇总}
            {ccccc}
            {\songtibf{假设} & \songtibf{路径} & \songtibf{间接效应} & \songtibf{BootCI} & \songtibf{结论} \\}
            {H2a & 理论型$\rightarrow$自主动机$\rightarrow$创业意愿 & 0.748 & [0.492, 1.031] & 完全中介 \\
            H2b & 实践型$\rightarrow$自主动机$\rightarrow$创业意愿 & 0.726 & [0.270, 1.682] & 完全中介 \\
            H2c & 融合型$\rightarrow$自主动机$\rightarrow$创业意愿 & 0.857 & [0.404, 1.615] & 完全中介 \\}
            {tab:mediation_summary}

\subsection{调节效应检验}

本研究采用分层回归的方法对外部动机的调节作用进行研究。在模型搭建时，把核心自变量和调节变量进行标准化处理，建立相关的交互作用项\nolinebreak[4]（自主动机$\times$外部动机）。因此，本研究对不同类型课程数据的调节效应做了详细的验证，具体的结果如表\ref{tab:moderation}所示。

由表\ref{tab:moderation}可知，在理论型课程样本中，加入控制变量、理论型课程、自主动机和外部动机后，自主动机与外部动机的交互项对创业意愿的回归系数为-0.148，对应的p值为0.332，未通过显著性检验。因此可以得出假设H3a没有得到足够的证据支持。

在实践型课程样本中，加入控制变量、实践型课程、自主动机和外部动机后，自主动机与外部动机的交互项对创业意愿的回归系数为-0.406，对应的p值为0.089，未通过显著性检验。因此可以得出假设H3b没有得到足够的证据支持。

在融合型课程样本中，加入控制变量、融合型课程、自主动机和外部动机后，自主动机与外部动机的交互项对创业意愿的回归系数为-0.535，对应的p值为0.048，小于0.05，通过了显著性检验。这表明外部动机对融合型课程中自主动机与创业意愿的关系具有显著的负向调节作用，即外部动机越强的学生，自主动机对创业意愿的正向影响越弱。因此，假设H3c成立。

综合分析结果可知，外部动机对自主动机与创业意愿关系的负向调节作用仅在融合型课程中显著\nolinebreak[4]（$\beta=-0.535$，$p=0.048$），在理论型课程\nolinebreak[4]（$\beta=-0.148$，$p=0.332$）和实践型课程\nolinebreak[4]（$\beta=-0.406$，$p=0.089$）中均不显著。因此，假设H3c成立，H3a和H3b不成立。

\tablemacro{htbp}
            {调节效应检验}
            {ccccccc}
            {\songtibf{假设} & \songtibf{课程类型} & \songtibf{交互项} & \songtibf{$\beta$} & \songtibf{p} & \songtibf{结论} \\}
            {H3a & 理论型 & 自主动机$\times$外部动机 & -0.148 & 0.332 & 不成立 \\
            H3b & 实践型 & 自主动机$\times$外部动机 & -0.406 & 0.089 & 不成立 \\
            H3c & 融合型 & 自主动机$\times$外部动机 & -0.535 & 0.048 & 成立 \\}
            {tab:moderation}

\noindent 注：表中报告的是标准化系数Beta；控制变量包括性别、年级、专业、创业经历、志愿服务。

H3a和H3b没有得到支持，可能的原因有以下几点：

第一，样本量限制。理论型课程样本68份、实践型53份、融合型仅33份。调节效应的检验需要较大的样本量才能达到足够的统计检验力，较小的样本量可能导致本应显著的调节效应无法被检测到。融合型课程的第二阶段调节效应显著\nolinebreak[4]（$p=0.045$），而其他不显著，可能与样本量差异有关。

第二，外部动机与课程的交互作用较弱。从相关分析来看，外部动机与三类课程的相关性均较低\nolinebreak[4]（理论型$r=0.289$，实践型$r=-0.04$，融合型$r=-0.25$），说明外部动机与课程感知的关联不强，这可能导致交互项的解释力有限。

第三，融合型课程的特殊性。只有融合型课程的调节效应显著，可能是因为融合型课程对学生的投入要求最高，外部动机的负面影响在此类课程中更容易显现。理论型课程中，学生即使动机不强也能被动接受知识；实践型课程中，实践操作本身具有一定的趣味性；而融合型课程需要学生在理论和实践两个层面都保持投入，外部动机的\nolinebreak[4]“挤出效应”在此类课程中最为明显。

\section{分析结果汇总}

本研究根据有效的问卷154份，对研究数据进行了详细的分析，同时用共同方法偏差检验来确定样本中不存在明显的共通性偏误。经过描述性统计分析之后，得出的结果表明所研究的对象特征也大致呈规律的分布。从相关性检验结果可知，被试创业意愿同个人内在驱动力以及各种课程模块都存在显著正相关关系。理论假设检验阶段，由多因素回归模型得出如下结论，所有的分类课程都会对个体的自主行为倾向起到正向促进的作用，H1a、H1b、H1c都成立，自主行为作为核心中介变量，它在整个类型课程和最终创业决策意图的关系链路上起到了完全中介的作用，H2a、H2b、H2c都成立。关于外部奖励机制影响机制的研究表明，在受试处于融合型课程的情境下，外部激励会对其学习参与度造成显著的抑制效应\nolinebreak[4]（H3c成立），但是在理论型和实践型课程的情境下，该因素并没有对认知过程或者决策判断产生显著的干扰\nolinebreak[4]（H3a、H3b不成立）。经过详细的研究发现如表\ref{tab:hypothesis_result}所示。

\tablemacro{htbp}
            {研究假设验证结果}
            {ccc}
            {\songtibf{假设序号} & \songtibf{假设内容} & \songtibf{验证结果} \\}
            {H1a & 理论型课程正向影响自主动机 & 成立 \\
            H1b & 实践型课程正向影响自主动机 & 成立 \\
            H1c & 融合型课程正向影响自主动机 & 成立 \\
            H2a & 自主动机中介理论型课程与创业意愿 & 成立 \\
            H2b & 自主动机中介实践型课程与创业意愿 & 成立 \\
            H2c & 自主动机中介融合型课程与创业意愿 & 成立 \\
            H3a & 外部动机负向调节理论型课程中自主动机与创业意愿 & 不成立 \\
            H3b & 外部动机负向调节实践型课程中自主动机与创业意愿 & 不成立 \\
            H3c & 外部动机负向调节融合型课程中自主动机与创业意愿 & 成立 \\}
            {tab:hypothesis_result}

\clearpage
